\section{Esito positivo: integrazione tramite Pull Request}
\label{sec:esito_positivo}

Quando la strategia passa la validazione walk-forward e rispetta i criteri di accettazione della issue, si procede con la sua integrazione nel sistema tramite pull request.

\subsection{Step 13a: Aggiornamento documentazione}

Prima di aprire la pull request, aggiornare i documenti che descrivono lo stato del sistema:

\begin{itemize}
  \item Creare (o aggiornare) il file \texttt{README.md} all'interno della cartella della strategia (\texttt{strategies/<nome\_strategia>/README.md}). Tale file deve contenere:
  \begin{itemize}
    \item Razionale economico della strategia;
    \item Riferimenti a letteratura;
    \item Parametri default e loro motivazione;
    \item Risultati OOS sintetici (Sharpe, CAGR, MaxDD, periodo coperto);
    \item Eventuali limiti noti.
  \end{itemize}
  \item Se la strategia introduce nuovi indicatori, aggiornare la sezione corrispondente di \texttt{GUIDE\_STRATEGIE\_CUSTOM.md};
  \item Se la strategia modifica numeri rilevanti dell'ensemble complessivo, aggiornare \texttt{PITCH\_FOF.md} con i nuovi valori.
\end{itemize}

\subsection{Step 14a: Pull request finale}

Aprire la pull request da \texttt{feature/<nome-strategia>} verso \texttt{main} seguendo il template descritto nella Sezione \ref{chap:processo_sviluppo}. La descrizione deve includere obbligatoriamente:

\begin{lstlisting}
Closes #<numero-issue>

## Strategia: <nome>

## Descrizione
<descrizione sintetica del razionale economico>

## Risultati walk-forward OOS
- Sharpe netto: X.XX (target era >Y.YY)
- CAGR netto: X.XX% (benchmark SPY: Y.YY%)
- Alpha vs SPY: +X.XX%
- MaxDD: -X.XX%
- Periodo coperto: YYYY-MM-DD -> YYYY-MM-DD

## Correlazione con stream esistenti
- vs momentum:  X.XX
- vs quality_lv: X.XX
- vs value:     X.XX
- vs bear_overlay: X.XX

## Test eseguiti
- [x] Walk-forward completo
- [x] Stress test 2008 / 2020 / 2022
- [x] Verifica decorrelazione vs fattori standard
- [x] Test unitari sull'implementation

## Note
<eventuali note: dipendenze nuove, exit policy scelta,
limiti di capacity osservati>
\end{lstlisting}

\subsection{Step 15a: Review e merge}

La pull request entra nel normale ciclo di review descritto nel Capitolo \ref{chap:processo_sviluppo}:

\begin{enumerate}
  \item Almeno un membro del team effettua la review;
  \item Particolare attenzione viene prestata a: presenza di lookahead bias, qualità dei test, ripetibilità del walk-forward, plausibilità dei numeri OOS;
  \item Eventuali modifiche richieste vengono implementate dall'autore;
  \item Una volta approvata, la PR viene mergiata con strategia \emph{Squash and merge};
  \item La issue collegata viene chiusa automaticamente.
\end{enumerate}

A questo punto la strategia è parte del sistema. Verrà discoverata automaticamente dal registry e potrà essere inclusa negli ensemble.

\subsection{Step 16a: Cleanup}

Dopo il merge:

\begin{enumerate}
  \item Cancellare il branch remoto da GitHub (l'interfaccia lo propone automaticamente);
  \item Cancellare il branch locale:
\begin{lstlisting}
git checkout main
git pull origin main
git branch -d feature/<nome-strategia>
\end{lstlisting}
  \item Verificare che la strategia appaia nella dashboard come stream selezionabile.
\end{enumerate}
